% Round 13. Owner: r13_analytic_descent.
% New internal results; published ingredients are identified explicitly.
\section{Exceptional Aeppli classes and actual primitive descent}
\label{r13d:section}

This section identifies the exact obstruction to actual primitive descent
through a smooth-center blowup as a quotient of exceptional Aeppli
classes.  It proves vanishing in a supplied retraction sector without
pushing an arbitrary actual form through a proper map.  It also exhibits a concrete
obstruction to extending the minimal finite trace to proper modifications,
even on its original scalar--curvature generators.  The obstruction is to
that pushforward operation, not to the existence of zero-curvature Green
forms.

Fix the original finite valued extension $K/\mathbf Q_p$.  All forms
below belong to Mihatsch's actual ordinary-topology sheaf $\mathcal B$;
$D=d'd''$.  The analytic assertions use Definition 4.2, Theorem 4.5,
Corollary 4.9 and Proposition 4.11 of \cite{M}.  The new deductions are
proved below.  No algebraic closure, changed absolute value, finite
dimensionality of Bott--Chern groups, or arbitrary proper actual-form
pushforward is assumed.

\subsection{Two elementary facts about actual forms}

\begin{lemma}[Restriction off a proper algebraic subset is faithful]
\label{r13d:dense-faithful}
Let $T/K$ be a smooth integral $n$-dimensional variety and let
$A\subsetneq T$ be a closed algebraic subset.  The restriction
\[
 \mathcal B^{p,q}(T^{\rm an})
       \longrightarrow\mathcal B^{p,q}((T\setminus A)^{\rm an})
\]
is injective.  Consequently pullback by a modification of smooth
integral varieties is injective on actual forms.
\end{lemma}

\begin{proof}
Suppose $\theta$ restricts to zero.  For each compact original-smooth
complementary test $\eta$, the actual top form $\theta\wedge\eta$ has
support contained in $A^{\rm an}$.  Corollary 4.9 of \cite{M} says that
the support of a top form avoids every algebraic subset of dimension
less than $n$.  Hence $\theta\wedge\eta=0$ and all current pairings
of $\theta$ vanish.  Proposition 4.11 then gives $\theta=0$ as an
actual form.  A modification is an isomorphism off a proper closed
subset; apply this assertion there to obtain the second statement.
\end{proof}

\begin{lemma}[An analytic retraction determines actual germs]
\label{r13d:retraction-germ}
Let $i:E\hookrightarrow Y$ be a closed immersion of pure-dimensional
boundaryless analytic spaces.  Suppose an ordinary open neighborhood
$N$ of $E$ admits an analytic retraction $\rho:N\to E$, so
$\rho i=\operatorname{id}_E$.  For every actual form $\beta$ on $N$
there is an ordinary open neighborhood $W\subset N$ of $E$ on which
\[
                         \beta=\rho^*i^*\beta.
\tag{R13.D1}\label{r13d:germ-equality}
\]
In particular, $i^*\beta=0$ implies that $\beta$ vanishes on an
ordinary neighborhood of $E$.  For this last implication it is enough
to have a covering of $E$ by such local retractions.
\end{lemma}

\begin{proof}
Near a point of $E$, choose a unit presentation
$\beta=t_f^\star\gamma$ with finitely many analytic units $f_j$ and
a Euclidean delta-form $\gamma$.  Shrink so that the units
$f_j^0=\rho^*i^*f_j$ are also defined.  The functions $f_j-f_j^0$
vanish on $E$.  Therefore the inequalities
\[
                          |f_j-f_j^0|<|f_j^0|
\]
define an ordinary open neighborhood of the part of $E$ under
consideration.  On it $|f_j|=|f_j^0|$ for every $j$.

Refine by the common tuple $(f,f^0)$ on any witnessing skeleton.
The two valuation maps are equal as linear maps on that same weighted
skeleton.  Their intrinsic tropical pullbacks of $\gamma$ therefore
agree.  This is equality of pullbacks on the same source; it is not
a restriction of an ambient Euclidean distribution to a nonlinear
subset.  Definition 4.2 and Lemma 4.4 of \cite{M} identify the actual
forms, and Theorem 4.5 identifies the second one with
$\rho^*i^*\beta$.  Taking the union of these neighborhoods proves
\eqref{r13d:germ-equality}.  The same proof on each member of a local
retraction covering proves the final assertion.
\end{proof}

This proof applies to arbitrary actual forms, not only to pure model
curvatures or to forms on a product family.  The retraction in the
hypothesis is an analytic morphism.  A continuous retraction onto a
Berkovich skeleton would not supply the pullbacks used here.

\begin{lemma}[Smooth closed immersions admit local analytic retractions]
\label{r13d:smooth-local-retraction}
Let $E\hookrightarrow Y$ be a closed immersion of smooth
$K$-varieties.  Every point of $E^{\rm an}$ has an ordinary analytic
neighborhood in $Y^{\rm an}$ carrying an analytic retraction onto
its intersection with $E^{\rm an}$.  Consequently an actual form
on $Y^{\rm an}$ whose restriction to $E^{\rm an}$ is zero vanishes
on an ordinary neighborhood of $E^{\rm an}$.
\end{lemma}

\begin{proof}
Choose a sufficiently small algebraic affine neighborhood $Y_0$
of the underlying scheme point, and write $E_0=E\cap Y_0$.
Smoothness supplies an etale map
$a:E_0\to\mathbf A_K^d$, with $d=\dim E_0$.
Its coordinate functions lift through the surjection
$\mathcal O(Y_0)\to\mathcal O(E_0)$, giving a map
$g:Y_0\to\mathbf A_K^d$ restricting to $a$.
This is the elementary coordinate construction used in Stacks,
Tags 054L and 0FUE \cite{Stacks}.

Form $W=Y_0\times_{\mathbf A_K^d}E_0$ and let
$\nu:W\to Y_0$ be the first projection.  It is etale, and there
is a diagonal lift $j:E_0\to W$ with $\nu j$ equal to the closed
immersion.  For any analytic point $e\in E_0^{\rm an}$, the map
of completed residue fields at $j(e)$ and $e$ is an isomorphism:
both closed immersions identify that field with $\mathscr H(e)$.
Berkovich's Theorem 3.4.1 \cite{BerEt} identifies finite etale
germs at $e$ with finite etale $\mathscr H(e)$-algebras.  Applied
to this degree-one branch, it gives ordinary open neighborhoods
on which $\nu^{\rm an}$ is an analytic isomorphism.

Shrink the target neighborhood $U$ further so that
$j(E_0^{\rm an}\cap U)$ is in the chosen inverse branch; this is
possible because $j$ is continuous and $E_0^{\rm an}$ is closed
in $Y_0^{\rm an}$.  The inverse branch followed by $W^{\rm an}\to
E_0^{\rm an}$ gives a morphism $r:U\to E_0^{\rm an}$ that is the
identity on $E_0^{\rm an}\cap U$.  Replace $U$ by the open
$U\cap (ir)^{-1}(U)$.  It is stable under $ir$, and $r$ is then
a retraction onto $E_0^{\rm an}\cap U$.  This construction is over
the original $K$ and applies also to nonrigid analytic points.
The last assertion follows from Lemma~\ref{r13d:retraction-germ}.
\end{proof}

\begin{lemma}[Degreewise extension from a smooth closed subvariety]
\label{r13d:restriction-surjective}
Let $E\hookrightarrow Y$ be a closed immersion of smooth projective
$K$-varieties.  Restriction
$i^*:\mathcal B^{a,b}(Y^{\rm an})\to\mathcal B^{a,b}(E^{\rm an})$
is surjective in every bidegree.  No extension commuting with
differentials is asserted.
\end{lemma}

\begin{proof}
Choose a finite ordinary neighborhood covering $U_j$ of the compact
space $E^{\rm an}$ carrying the retractions in the preceding lemma.
Add $Y^{\rm an}\setminus E^{\rm an}$ to obtain an open covering of
$Y^{\rm an}$.  Smooth partitions of unity subordinate to such
ordinary covers exist by Chambert-Loir--Ducros, Proposition 3.3.6
of the original 2012 version \cite{CLDold}.
The smooth coefficients multiply actual forms by Remark 4.3(1)
of \cite{M}.

For $\theta\in\mathcal B^{a,b}(E^{\rm an})$, multiply each local
pullback $r_j^*\theta$ by its corresponding smooth cutoff and
extend the product by zero outside $U_j$.  The supports of the
cutoffs lie inside $U_j$, so these zero extensions are actual
forms.  Their finite sum restricts to $\theta$, because the
cutoff associated to $Y\setminus E$ restricts to zero and the
remaining cutoffs sum to one on $E$.  This proves surjectivity.
The derivatives of the cutoffs explain why this construction is
only degreewise and is not claimed to preserve closedness.
\end{proof}

\subsection{Relative forms and an exceptional extension criterion}

Let $q:Y\to X$ be a projective birational morphism between smooth
integral projective $K$-varieties.  Fix a smooth closed subvariety
$s:Z\hookrightarrow X$ such that $q$ is an isomorphism over
$X\setminus Z$.  Suppose the inverse image of $Z$ has underlying
closed analytic space a smooth closed subvariety $i:E\hookrightarrow Y$.
Write $q_E:E\to Z$ for the induced morphism.  By
Lemma~\ref{r13d:smooth-local-retraction} the following property holds:
\[
 \begin{gathered}
 \text{the immersion }E^{\rm an}\hookrightarrow Y^{\rm an}
 \\
 \text{ admits local analytic retractions on an ordinary}
 \\
 \text{ neighborhood covering of }E^{\rm an}.
 \end{gathered}
\tag{R13.D2}\label{r13d:local-retractions}
\]
In particular this setup applies to every blowup of a smooth center
in a smooth projective variety: its exceptional divisor is a smooth
projective bundle over the center by Stacks, Tag 0FUT.
The displayed local retractions will be used only to identify germs;
global extension of a $D$-closed form remains a separate problem.

Define the relative double complexes
\[
 \mathcal R_X^{p,q}=
       \ker\bigl(s^*:\mathcal B^{p,q}(X^{\rm an})
                              \to\mathcal B^{p,q}(Z^{\rm an})\bigr),
 \qquad
 \mathcal R_Y^{p,q}=
       \ker\bigl(i^*:\mathcal B^{p,q}(Y^{\rm an})
                              \to\mathcal B^{p,q}(E^{\rm an})\bigr).
\]

\begin{proposition}[Relative actual descent]
\label{r13d:relative-isomorphism}
In the stated smooth modification setup, pullback induces an isomorphism
of bigraded differential algebras
\[
                          q^*:\mathcal R_X\xrightarrow{\ \sim\ }
                                      \mathcal R_Y.
\tag{R13.D3}\label{r13d:relative-isom}
\]
Its inverse $T$ is the following concrete operation: transport a
relative form through the isomorphism off the exceptional locus,
and extend it by zero over an ordinary neighborhood of $Z$.
It commutes with $d'$, $d''$, and $D$.
\end{proposition}

\begin{proof}
Pullback maps $\mathcal R_X$ into $\mathcal R_Y$ by the commutative
square of the two closed immersions, and is injective by
Lemma~\ref{r13d:dense-faithful}.  Let $\lambda\in\mathcal R_Y$.
Lemma~\ref{r13d:retraction-germ} supplies an ordinary open
$W\subset Y^{\rm an}$ containing $E^{\rm an}$ on which
$\lambda=0$.

Set
\[
                 V=X^{\rm an}\setminus
                        q^{\rm an}(Y^{\rm an}\setminus W).
\tag{R13.D4}\label{r13d:saturated-neighborhood}
\]
The set removed is closed: $Y^{\rm an}$ is compact, $X^{\rm an}$ is
Hausdorff, and the map is continuous.  Thus $V$ is an ordinary open
containing $Z^{\rm an}$, and $(q^{\rm an})^{-1}V\subset W$.
Transport $\lambda$ to $(X\setminus Z)^{\rm an}$ and glue it with
zero on $V$.  The forms agree on the overlap because $\lambda=0$
on $q^{-1}V$.  The result $T\lambda$ is an actual form by the
ordinary sheaf property, and its pullback is exactly $\lambda$.
It vanishes near $Z$ and thus belongs to $\mathcal R_X$.

The operator is independent of $W$ and $V$ by uniqueness off $Z$
and Lemma~\ref{r13d:dense-faithful}.  Its compatibility with products
and the differential operators can be checked on the displayed open
cover, where it is either transport through an isomorphism or zero
extension.  This proves all assertions.
\end{proof}

\begin{theorem}[A short exact double complex for a smooth modification]
\label{r13d:double-complex}
In the preceding setup there is a short exact sequence of bigraded
vector spaces, compatible with both $d'$ and $d''$,
\[
0\longrightarrow\mathcal B(X^{\rm an})
 \xrightarrow{\ (q^*,s^*)\ }
 \mathcal B(Y^{\rm an})\oplus\mathcal B(Z^{\rm an})
 \xrightarrow{\ i^*-q_E^*\ }
 \mathcal B(E^{\rm an})\longrightarrow0.
\tag{R13.D14}\label{r13d:short-exact}
\]
Its exactness is degreewise.  No splitting as a double complex is
asserted here.
\end{theorem}

\begin{proof}
Injectivity on the left is Lemma~\ref{r13d:dense-faithful}.
Surjectivity on the right follows from
Lemma~\ref{r13d:restriction-surjective} applied to $i$.
Suppose a pair $(\beta,\zeta)$ satisfies $i^*\beta=q_E^*\zeta$.
Extend $\zeta$ to an actual form $\eta$ on $X$ by the same lemma
applied to $s$.  Then $\beta-q^*\eta$ is relative on $Y$, so
\[
                         \alpha=\eta+T(\beta-q^*\eta)
\]
has $q^*\alpha=\beta$ and $s^*\alpha=\zeta$.
This proves middle exactness.  Compatibility with the differentials
holds for all displayed arrows because they are analytic pullbacks.
\end{proof}

\begin{theorem}[Extend the exceptional restriction, then descend]
\label{r13d:closed-extension}
In the preceding smooth modification setup, let
$\Omega\in\mathcal B^{r,r}(X^{\rm an})$ and suppose an actual form
$\beta\in\mathcal B^{r-1,r-1}(Y^{\rm an})$ satisfies
\[
                             D\beta=q^*\Omega.
\]
If there exists a global actual form $\chi$ of the same bidegree on
$Y^{\rm an}$ with
\[
                         D\chi=0,\qquad i^*\chi=i^*\beta,
\tag{R13.D5}\label{r13d:closed-extension-condition}
\]
then $\Omega$ has the global actual primitive
\[
                           \beta_X=T(\beta-\chi).
\tag{R13.D6}\label{r13d:extended-primitive}
\]
\end{theorem}

\begin{proof}
The difference lies in $\mathcal R_Y$, so the preceding proposition
defines $\beta_X$ and gives $q^*\beta_X=\beta-\chi$.  Apply $D$:
\[
                      q^*(D\beta_X-\Omega)=0.
\]
Faithfulness of actual pullback by a modification proves
$D\beta_X=\Omega$.
\end{proof}

The extension requirement concerns $D$-closed forms; separate
$d'$- and $d''$-closedness of $\chi$ is not needed.  It is a
substantive additional condition.  In particular, an arbitrary
primitive supplied on an alteration is not asserted to satisfy it.

\begin{corollary}[A global exceptional retraction]
\label{r13d:global-retraction}
Suppose that $i:E\hookrightarrow Y$ has a global analytic retraction
$\rho:Y^{\rm an}\to E^{\rm an}$.  If
$q_E^*s^*\Omega=0$ and $D\beta=q^*\Omega$, then
\[
          \boxed{\quad
          \beta_X=T\bigl(\beta-\rho^*i^*\beta\bigr),
          \qquad D\beta_X=\Omega.\quad}
\tag{R13.D7}\label{r13d:retracted-primitive}
\]
In particular, for $r>\dim Z$ the pullback
\[
 q^*:H_{\mathrm{BC,act}}^{r,r}(X)
                       \longrightarrow H_{\mathrm{BC,act}}^{r,r}(Y)
\]
is injective.
\end{corollary}

\begin{proof}
Take $\chi=\rho^*i^*\beta$.  Since pullback commutes with $D$,
\[
             D\chi=\rho^*i^*q^*\Omega
                        =\rho^*q_E^*s^*\Omega=0.
\]
Also $i^*\chi=i^*\beta$.  Theorem~\ref{r13d:closed-extension}
gives the formula.  If $r>\dim Z$, the space
$\mathcal B^{r,r}(Z^{\rm an})$ is zero by Corollary 4.9 of \cite{M},
so the additional vanishing holds automatically.  Applying the formula
to every separately closed $\Omega$ whose pullback is $D$-exact
proves the Bott--Chern assertion.
\end{proof}

\subsection{The complete exceptional Aeppli obstruction}

For a smooth projective $T/K$, define the algebraic Aeppli quotient
\[
 H_{\mathrm A,\mathrm{act}}^{k,k}(T)
 =\frac{\ker\bigl(D:\mathcal B^{k,k}(T^{\rm an})
                          \to\mathcal B^{k+1,k+1}(T^{\rm an})\bigr)}
 {d'\mathcal B^{k-1,k}(T^{\rm an})+
          d''\mathcal B^{k,k-1}(T^{\rm an})}.
\tag{R13.D15}\label{r13d:aeppli}
\]
Spaces in negative bidegrees are zero.  The denominator is contained
in the numerator by the differential identities.  All quotients in
this subsection are plain vector-space quotients; no closed-range
or finite-dimensionality assertion is needed.

\begin{theorem}[The Aeppli--Bott--Chern connecting sequence]
\label{r13d:connecting}
For every $k\geq0$, the following sequence is exact at its two
middle terms:
\[
\begin{split}
 H_{\mathrm A,\mathrm{act}}^{k,k}(Y)\oplus
 H_{\mathrm A,\mathrm{act}}^{k,k}(Z)
 &\xrightarrow{\ i^*-q_E^*\ }
 H_{\mathrm A,\mathrm{act}}^{k,k}(E)
 \xrightarrow{\ \partial_D\ }
 H_{\mathrm{BC,act}}^{k+1,k+1}(X)\\
 &\xrightarrow{\ (q^*,s^*)\ }
 H_{\mathrm{BC,act}}^{k+1,k+1}(Y)\oplus
 H_{\mathrm{BC,act}}^{k+1,k+1}(Z).
\end{split}
\tag{R13.D16}\label{r13d:connecting-sequence}
\]
An explicit connecting map is
\[
               \partial_D[\theta]=[T(D\widetilde\theta)],
       \qquad i^*\widetilde\theta=\theta,\quad D\theta=0,
\tag{R13.D17}\label{r13d:connecting-formula}
\]
where $\widetilde\theta$ is any degreewise actual extension to $Y$.
\end{theorem}

\begin{proof}
The extension exists by Lemma~\ref{r13d:restriction-surjective}.
Since $i^*D\widetilde\theta=D\theta=0$, the form
$D\widetilde\theta$ lies in $\mathcal R_Y$.  Its descent
$\Omega=T(D\widetilde\theta)$ is separately closed, because $T$
commutes with both differentials, and it restricts to zero on $Z$.
Thus the displayed Bott--Chern class is defined.

If the extension is changed, the difference is relative and hence
equals $q^*\lambda$ for a global actual $\lambda$ on $X$.
The descended curvature changes by $D\lambda$, so its class is
unchanged.  If $\theta$ changes by $d'a+d''b$, extend $a,b$
degreewise to $Y$ and change the extension by
$d'\widetilde a+d''\widetilde b$.  Its $D$ is zero, so the formula
passes to the Aeppli quotient.

We check exactness at the Aeppli term.  The image of a $D$-closed
form on $Y$ maps to zero, by using that same form as extension.
If $\theta=q_E^*\zeta$ with $D\zeta=0$ on $Z$, extend $\zeta$
to some $\eta$ on $X$ and use $q^*\eta$ as extension on $Y$.
Its derivative descends to $D\eta$, because
$s^*D\eta=D\zeta=0$; its Bott--Chern class is zero.  These facts
show that the preceding image is contained in the kernel.

Conversely suppose $\partial_D[\theta]=0$.  With the extension
above write $T(D\widetilde\theta)=D\eta$ for an actual $\eta$
on $X$.  Then
$\chi=\widetilde\theta-q^*\eta$ is $D$-closed on $Y$, and
$\zeta=s^*\eta$ is $D$-closed on $Z$, since the descended curvature
restricts to zero.  The identity
\[
                            \theta=i^*\chi+q_E^*\zeta
\]
proves membership in the preceding image (the sign is immaterial
for its image as a vector subspace).

Finally, the image of $\partial_D$ has zero pullback Bott--Chern
class on $Y$ and zero restriction on $Z$.  Conversely let a
separately closed actual $\Omega$ represent a class in the kernel
of $(q^*,s^*)$.  Choose primitives
\[
                        q^*\Omega=D\beta,\qquad
                        s^*\Omega=D\zeta.
\]
Extend $\zeta$ degreewise to $\eta$ on $X$.  Then
$\Omega_0=\Omega-D\eta$ restricts literally to zero on $Z$, and
$\beta_0=\beta-q^*\eta$ satisfies $D\beta_0=q^*\Omega_0$.
The exceptional restriction $\theta=i^*\beta_0$ is $D$-closed.
Using $\beta_0$ in \eqref{r13d:connecting-formula} gives
$\partial_D[\theta]=[\Omega_0]=[\Omega]$.
This proves exactness at the Bott--Chern term by direct
representative calculations, without an assumed long exact
sequence for Bott--Chern cohomology.
\end{proof}

\begin{corollary}[An exact criterion for the curve-degree descent kernel]
\label{r13d:aeppli-kernel}
For $r>\dim Z$, the connecting map gives a canonical isomorphism
\[
 \boxed{\quad
 \frac{H_{\mathrm A,\mathrm{act}}^{r-1,r-1}(E)}
 {i^*H_{\mathrm A,\mathrm{act}}^{r-1,r-1}(Y)
        +q_E^*H_{\mathrm A,\mathrm{act}}^{r-1,r-1}(Z)}
 \ \simeq\ 
 \ker\bigl(q^*:H_{\mathrm{BC,act}}^{r,r}(X)
                         \to H_{\mathrm{BC,act}}^{r,r}(Y)\bigr).
 \quad}
\tag{R13.D18}\label{r13d:aeppli-kernel-iso}
\]
In particular, this applies to $r=n-1$ for every blowup of a
nonempty smooth center of codimension at least two in an $n$-fold.
For a specific model-curvature class, only its corresponding
exceptional quotient class must vanish; the whole quotient need
not vanish.
\end{corollary}

\begin{proof}
The rightmost Bott--Chern group of $Z$ in
\eqref{r13d:connecting-sequence} is zero when $r>\dim Z$.
The asserted isomorphism is therefore precisely the two exactness
statements of the theorem.  In the stated curve degree,
$\dim Z\leq n-2<n-1$.
\end{proof}

Here is the constructive criterion for an individual curvature.
Let $r>\dim Z$, let $\Omega$ be separately closed on $X$, and
suppose $D\beta=q^*\Omega$.  Its exceptional class is
\[
 \varepsilon_q(\Omega)
     =[i^*\beta]\quad\text{in the quotient on the left of
                          \eqref{r13d:aeppli-kernel-iso}}.
\tag{R13.D19}\label{r13d:individual-obstruction}
\]
It is independent of the primitive $\beta$, because two primitives
differ by a global $D$-closed form on $Y$.  Changing the representative
$\Omega$ by a global $D$ changes $i^*\beta$ by the pullback of a
$D$-closed form on $Z$, so the class depends only on $[\Omega]$.
It vanishes if and only if $\Omega$ has a global actual primitive
on $X$.

More explicitly, a vanishing certificate consists of $D$-closed
actual forms $\chi$ on $Y$ and $\zeta$ on $Z$, and actual forms
$a,b$ on $E$, such that
\[
                    i^*\beta=i^*\chi+q_E^*\zeta+d'a+d''b.
\]
Extend $a,b$ to $Y$ and absorb $d'\widetilde a+d''\widetilde b$
into $\chi$.  Extend $\zeta$ to any actual $\eta$ on $X$.
Then the formula
\[
                  \boxed{\quad
                  \Psi_X=\eta+T(\beta-q^*\eta-\chi),
                  \qquad D\Psi_X=\Omega\quad}
\tag{R13.D20}\label{r13d:certificate-primitive}
\]
is valid.  The input of $T$ has zero exceptional restriction, and
its derivative is $q^*(\Omega-D\eta)$, which proves the equation.
This is an equivalence for primitive descent.  It is not a claimed
equivalence between strong G and vanishing of the R11 sufficient
model-curvature obstruction.

\subsection{The blowup of a zero section in a split projective bundle}

Let $S/K$ be smooth integral projective of dimension $d$, and let
$V$ be a vector bundle of rank $c\geq2$.  Use the projective-bundle
convention of lines.  Put
\[
 X=\mathbf P_S(\mathcal O_S\oplus V),\qquad
 Z=\mathbf P_S(\mathcal O_S)\simeq S,
 \qquad
 E=\mathbf P_S(V).
\]
Write $p:X\to S$, $\pi:E\to S$, and let
$L=\mathcal O_E(-1)\subset\pi^*V$ be the tautological line.  There
is an isomorphism
\[
       Y:=\operatorname{Bl}_{Z}X
          \simeq\mathbf P_E(\mathcal O_E\oplus L).
\tag{R13.D8}\label{r13d:bundle-blowup}
\]
Its projection $\rho:Y\to E$ retracts the exceptional section
$i:E=\mathbf P_E(\mathcal O_E)\hookrightarrow Y$.  The blowdown
$q:Y\to X$ is induced by the inclusion
$\mathcal O_E\oplus L\subset\pi^*(\mathcal O_S\oplus V)$, and
\[
                            pq=\pi\rho.
\tag{R13.D9}\label{r13d:bundle-square}
\]

\begin{proof}[Verification of \eqref{r13d:bundle-blowup}]
Trivialize $V$ on an open of $S$ and write a point of $X$ as
$[a:v_1:\cdots:v_c]$.  The indicated $Y$ parametrizes a line
$[\ell_1:\cdots:\ell_c]\in\mathbf P^{c-1}$ together with a point
$[a:b\ell_1:\cdots:b\ell_c]$ of the corresponding projective line.
The map to $X$ forgets the direction $\ell$ and is an isomorphism
where $(v_1,\ldots,v_c)\ne0$.

On $a\ne0$, its equations are the standard incidence equations
$v_i\ell_j=v_j\ell_i$ for the blowup of the coordinate ideal
$(v_1,\ldots,v_c)$ in affine $c$-space.  On $a=0$ the vector
$(v_1,\ldots,v_c)$ is nonzero, so it is an isomorphism there.  These
descriptions cover $X$ and are compatible with changes of frame of
$V$, proving the global identification by the usual affine charts
of the Rees-algebra blowup.  The locus $b=0$ is the exceptional
section, and the projection to $E$ restricts to its identity.
The identities claimed above follow from these descriptions.
\end{proof}

\begin{theorem}[A split blowup reduces injectivity to its base]
\label{r13d:split-kernel}
In \eqref{r13d:bundle-blowup}, for every diagonal degree $r$, the
kernel of
\[
                  q^*:H_{\mathrm{BC,act}}^{r,r}(X)
                                  \to H_{\mathrm{BC,act}}^{r,r}(Y)
\]
is identified by restriction to $Z\simeq S$ with
\[
                  \ker\bigl(\pi^*:H_{\mathrm{BC,act}}^{r,r}(S)
                                    \to H_{\mathrm{BC,act}}^{r,r}(E)\bigr).
\tag{R13.D10}\label{r13d:kernel-base}
\]
Thus $q^*$ is injective for $r>d$.  If $V$ has a line subbundle
with locally free quotient, so that $\pi$ has a section
$\sigma:S\to E$, then $q^*$ is injective in every degree.
\end{theorem}

\begin{proof}
The two genuine retractions split the double complexes:
\[
 \mathcal B(X^{\rm an})\simeq\mathcal R_X\oplus\mathcal B(S^{\rm an}),
 \qquad
 \alpha\longmapsto(\alpha-p^*s^*\alpha,s^*\alpha),
\]
\[
 \mathcal B(Y^{\rm an})\simeq\mathcal R_Y\oplus\mathcal B(E^{\rm an}),
 \qquad
 \beta\longmapsto(\beta-\rho^*i^*\beta,i^*\beta).
\]
Each is a direct-sum decomposition of bigraded vector spaces compatible
with both differentials.  We do not claim it is a direct product of
algebras.  By \eqref{r13d:bundle-square}, pullback $q^*$ becomes
the block-diagonal map $(q^*|_{\mathcal R_X},\pi^*)$.
Proposition~\ref{r13d:relative-isomorphism} makes its first block
an isomorphism.  Bott--Chern cohomology commutes with this finite
direct sum purely algebraically, giving
\eqref{r13d:kernel-base}.  In degree $r>d$ its source on $S$ is
zero.  If a section $\sigma$ exists, $\sigma^*\pi^*=\operatorname{id}$
on actual forms and hence on Bott--Chern groups, proving the final
assertion.
\end{proof}

When the section $\sigma$ is supplied, an explicit linear left inverse
of $q^*$ on the entire actual double complex is
\[
 \mathcal L(\beta)=
    T(\beta-\rho^*i^*\beta)+p^*\sigma^*i^*\beta.
\tag{R13.D11}\label{r13d:left-inverse}
\]
It commutes with $d'$ and $d''$.  Substituting $q^*\alpha$ gives
$\alpha-p^*s^*\alpha+p^*s^*\alpha=\alpha$.  The operator is not
asserted to be multiplicative, and it is not the current pushforward
along $q$.

\begin{corollary}[Consequences for the model-curvature obstruction]
\label{r13d:strong-G}
Let $a\in\operatorname{CH}^r(X)_{\mathbf Q}$ in the projective bundle
setting.  If
\[
                         \mathfrak o_Y^r(q^*a)=0,
\]
and either $r>d$ or $V$ has the line subbundle specified above, then
$\mathfrak o_X^r(a)=0$.  Every prescribed representative of $a$
therefore has a strong zero-curvature actual Green form over the
original $K$, on its exact cycle complement and on every compact
actual test.

In particular, let $n=d+c\geq3$ and let $a$ be geometrically
homologically trivial of codimension $n-1$.  A supplied regular
projective model of $Y$ is sufficient for this conclusion; no regular
model of $X$ is part of this implication's hypotheses.
\end{corollary}

\begin{proof}
The R11 curvature map is compatible with actual pullback, so its
value on $q^*a$ is $q^*\mathfrak o_X^r(a)$.  Apply
Theorem~\ref{r13d:split-kernel}.  The exact-support Green-pair
construction in R11 then subtracts the obtained global actual
primitive from the pair attached to the prescribed cycle.  The
identity $D\beta_X=\Omega$ is an identity of actual forms, and
Stokes applies to every compact actual test.

For the final assertion, $n-1=d+c-1>d$ because $c\geq2$.
The pullback of a geometrically homologically trivial class remains
so.  R12's curve theorem on the supplied regular model of $Y$ gives
$\mathfrak o_Y^{n-1}(q^*a)=0$ and the first assertion applies.
This is a conditional extension of the regular-model criterion to
this explicit non-finite generic modification.  It does not establish
the existence of such a regular model for every $Y$.
\end{proof}

\subsection{Why a universal proper trace on the old generators fails}

\begin{proposition}[The actual range of birational current pushforward]
\label{r13d:range}
Let $q:Y\to X$ be a projective birational morphism between smooth
integral projective $n$-folds, and let
$\beta\in\mathcal B^{a,b}(Y^{\rm an})$.  Its proper CLD current pushforward
is represented by an actual form $\alpha$ on $X^{\rm an}$ if and
only if
\[
                                 \beta=q^*\alpha.
\tag{R13.D12}\label{r13d:actual-range}
\]
Equivalently, a birational current pushforward produces an actual
form precisely for forms that already descend by actual pullback.
\end{proposition}

\begin{proof}
If $q_*[\beta]=[\alpha]$, restrict to an algebraic open on which
$q$ is an isomorphism.  Equality of currents there and Proposition
4.11 of \cite{M} imply $\beta=q^*\alpha$ on its inverse image.
Lemma~\ref{r13d:dense-faithful} extends this equality to $Y$.

Conversely suppose $\beta=q^*\alpha$.  For a compact smooth
complementary test $\eta$, the top forms $\alpha\eta$ and
$q^*(\alpha\eta)$ avoid the exceptional algebraic sets by
Corollary 4.9.  Their integrals can therefore be computed on the
isomorphism locus, where they agree.  Thus
\[
             \int_Y\beta\,q^*\eta=\int_X\alpha\eta,
\]
which is the required equality of currents.  The same argument is
valid for every compact actual complementary test.
\end{proof}

\begin{theorem}[A proper trace fails even on pure PL curvatures]
\label{r13d:counterexample}
For every $n\geq2$ and every $1\leq k\leq n-1$, there are global
PL scalars $v_1,\ldots,v_k$ on
$Y=\operatorname{Bl}_{p}\mathbf P^n_K$, with $p\in\mathbf P^n(K)$,
such that
\[
                      \beta=Dv_1\cdots Dv_k
                           \in\mathscr M^k(Y^{\rm an})
\]
has no actual form representing $q_*[\beta]$ on $\mathbf P^{n,\rm an}$.
Here $\mathscr M^k$ is exactly the scalar--curvature image sheaf
used in the R11 minimal finite trace.  In addition, $D\beta=0$ and
$\beta$ itself is globally $D$-exact.
\end{theorem}

\begin{proof}
Use \eqref{r13d:bundle-blowup} over $S=\operatorname{Spec}K$ with
$V=K^n$.  Then $E=\mathbf P^{n-1}$, and
$\rho:Y\to E$ retracts the exceptional divisor.  Write $m=n-1$.
On the standard dense torus of $E$, put
\[
 t_j=v(x_j/x_0),\qquad R(t)=\max_{1\leq j\leq m}|t_j|,
 \qquad b(t)=\max(0,2a-R(t)),
\]
where $a=v(\varpi)>0$ for a uniformizer of the original $K$.  Define
\[
                        u_j(t)=\min(|t_j|,b(t)).
\tag{R13.D13}\label{r13d:toric-potentials}
\]
These are finite PL functions.  They vanish for $R(t)\geq2a$.
The inverse image of $R(t)\leq2a$ is a compact affinoid polyannulus
contained in the torus, hence closed in $E^{\rm an}$.  Extending
$u_j$ by zero off the torus consequently gives a global PL scalar
on $E^{\rm an}$: the torus and the complement of that compact
polyannulus are an ordinary open cover, with equal functions on
the overlap.

On the open cube $R(t)<a$, one has $u_j=|t_j|$.  Thus the actual
form
\[
                              C_k=Du_1\cdots Du_k
\]
restricts on the standard weighted torus skeleton there to the
nonzero tropical intersection
\[
                         2^k[t_1=\cdots=t_k=0].
\]
The coordinate hyperplanes meet transversally, and each corner
locus of $|t_j|$ has weight two.  Hence $C_k\ne0$ as an actual
form by its skeletal presentation.

Set $v_j=\rho^*u_j$ and $\beta=\rho^*C_k$.
Then $i^*\beta=C_k\ne0$.  If its current pushforward were
represented by an actual $\alpha\in\mathcal B^{k,k}(\mathbf P^{n,\rm an})$,
Proposition~\ref{r13d:range} would give $\beta=q^*\alpha$.
Every positive-degree actual form vanishes on an ordinary
neighborhood of the rigid point $p$, by Corollary 4.9 of \cite{M}.
Its pullback therefore vanishes near the whole exceptional divisor,
contradicting $i^*\beta=C_k\ne0$.

Each curvature factor is separately closed, so $D\beta=0$ and
\[
                         \beta=D(v_1Dv_2\cdots Dv_k).
\]
This also proves the final assertion.
\end{proof}

For $n=3$ and $k=1$, this is already the bidegree of a curve-cycle
Green primitive.  It prevents a universal extension of the R11
finite trace even if its domain is kept to pure products of
curvatures of global PL functions.  Conversely, in this example
the retraction correction in \eqref{r13d:retracted-primitive}
subtracts the entire unwanted form: $\rho^*i^*\beta=\beta$.
The example therefore supports the need to adjust a primitive by
a $D$-closed exceptional contribution; it supplies no counterexample
to strong G or to injectivity of the curvature obstruction.

\paragraph{Remaining scope.}
The new general input is the exact exceptional Aeppli quotient
\eqref{r13d:aeppli-kernel-iso} for smooth-center blowups, together
with its actual primitive formula.  Vanishing of that quotient is
proved in the explicit retraction sector, not for all blowups.
A general alteration does not automatically supply the needed
exceptional vanishing certificate.  The projective-bundle calculation
does not assert a general non-archimedean Bott--Chern blowup formula,
nor does it remove the regular-model condition for all curve cycles.
All positive conclusions are identities of actual forms over the
original field and therefore preserve the exact-support and all-test
requirements when composed with the inherited Green-pair construction.
